% Options for packages loaded elsewhere
\PassOptionsToPackage{unicode}{hyperref}
\PassOptionsToPackage{hyphens}{url}
%
\documentclass[
  10pt,
]{article}
\usepackage{amsmath,amssymb}
\usepackage{iftex}
\ifPDFTeX
  \usepackage[T1]{fontenc}
  \usepackage[utf8]{inputenc}
  \usepackage{textcomp} % provide euro and other symbols
\else % if luatex or xetex
  \usepackage{unicode-math} % this also loads fontspec
  \defaultfontfeatures{Scale=MatchLowercase}
  \defaultfontfeatures[\rmfamily]{Ligatures=TeX,Scale=1}
\fi
\usepackage{lmodern}
\ifPDFTeX\else
  % xetex/luatex font selection
\fi
% Use upquote if available, for straight quotes in verbatim environments
\IfFileExists{upquote.sty}{\usepackage{upquote}}{}
\IfFileExists{microtype.sty}{% use microtype if available
  \usepackage[]{microtype}
  \UseMicrotypeSet[protrusion]{basicmath} % disable protrusion for tt fonts
}{}
\makeatletter
\@ifundefined{KOMAClassName}{% if non-KOMA class
  \IfFileExists{parskip.sty}{%
    \usepackage{parskip}
  }{% else
    \setlength{\parindent}{0pt}
    \setlength{\parskip}{6pt plus 2pt minus 1pt}}
}{% if KOMA class
  \KOMAoptions{parskip=half}}
\makeatother
\usepackage{xcolor}
\usepackage[margin=0.85in]{geometry}
\setlength{\emergencystretch}{3em} % prevent overfull lines
\providecommand{\tightlist}{%
  \setlength{\itemsep}{0pt}\setlength{\parskip}{0pt}}
\setcounter{secnumdepth}{-\maxdimen} % remove section numbering
\usepackage{amsmath,amssymb,amsthm}
\usepackage{longtable}
\usepackage{mathrsfs}
\ifLuaTeX
  \usepackage{selnolig}  % disable illegal ligatures
\fi
\usepackage{bookmark}
\IfFileExists{xurl.sty}{\usepackage{xurl}}{} % add URL line breaks if available
\urlstyle{same}
\hypersetup{
  pdftitle={Fixed references, fiber constancy, and finite-data descent for P16.1},
  pdfauthor={Primary-source reconstruction; unverified mathematical draft},
  hidelinks,
  pdfcreator={LaTeX via pandoc}}

\title{Fixed references, fiber constancy, and finite-data descent for
P16.1}
\author{Primary-source reconstruction; unverified mathematical draft}
\date{5 October 2026}

\begin{document}
\maketitle

\subsection{1. Exact disposition and
hypotheses}\label{exact-disposition-and-hypotheses}

The original P16.1 is \textbf{not solved}. The result below is a proof
draft of fiber constancy and a conditional reduction of the remaining
equation theorem to a precisely specified countable original-class
membership gate. Candidate=false. All axes of the original family and
the completion criterion are retained byte for byte in the preserved
input and task file. Nothing about the independent countable Runge work
is assumed or inspected.

The equation domain considered here is the domain specified in the
assignment: \(k\) is an arbitrary algebraically closed,
\textbf{trivially valued}, characteristic-zero field; \(X\) is a reduced
irreducible geometrically unibranch projective \(k\)-variety; and \(L\)
is an actual ample rational line bundle. Clearing a positive tensor
denominator reduces to honest line bundles when necessary. All
dimensions are allowed. For dimension zero, \(X=\operatorname{Spec}k\),
the empty curvature product has raw mass one, and the equation is
immediate. The arguments involving curves and nonempty products below
concern positive dimension. We do not remove projectivity, ampleness,
characteristic zero, algebraizability, or the trivial valuation. In
particular no arbitrary-field finite-energy theorem follows from the
auxiliary nontrivially valued fields used for fiber tests.

Let \(\pi:Y\to X\) be a projective birational morphism with \(Y\) smooth
irreducible projective, factoring through the finite normalization
\(X^\nu\). Such a resolution is obtained in characteristic zero; when
the task-granted contact bridge is invoked, its particular
compactification, projective birational domination with smooth
projective source, and resolution-witness antecedents remain binding. A
``smooth projective morphism'' is not inferred. The accepted partial
argument is mathematical background only in its exact contact scope. It
does not assert arbitrary weak approximation, singular EP, onto descent,
or the countable main theorem.

The cited primary assertions are indexed by nonnegative integers in
\texttt{premises.json}. Source units are literal UTF-8 decoded text,
with inclusive one-based LF coordinates computed using
\texttt{split(\textquotesingle{}\textbackslash{}n\textquotesingle{})};
source provenance distinguishes supplied PDF/text files from
web-rendered snapshots. External source theorems are inputs, not new
proofs of those theorems. Their applications and the elementary
arguments below are unverified drafts.

\subsection{2. A numerical specialization
lemma}\label{a-numerical-specialization-lemma}

We first prove the numerical fact needed for the \textbf{varying-class}
net. Let \(Z\) be projective over algebraically closed \(k\), let
\(\Omega/k\) be an algebraically closed extension, and let \(D\) be a
line bundle numerically trivial on \(Z\). Then \(D_\Omega\) is
numerically trivial on \(Z_\Omega\). More generally, for a fixed proper
integral curve \(T/\Omega\) and morphism \(a:T\to Z_\Omega\), the rule
\[[D]\longmapsto \deg_\Omega a^*D_\Omega                                      \tag{2.1}\]
factors through a finite continuous linear functional on
\(N^1(Z)_\mathbb R\).

Here is the finite-data proof. Embed \(T\) in a projective space and
write its equations, the finite morphism charts of \(a\), and the
equations identifying its generic fiber over a finitely generated
integral \(k\)-subalgebra \(A\subset\Omega\). Clearing the finite
denominators gives a projective finitely presented family
\(\mathcal T\to\operatorname{Spec}A\) and a morphism
\(\mathcal T\to Z_A\), with \(\mathcal T_\Omega=T\). Shrink this
integral base once so that \(\mathcal T\) is flat. This is generic
flatness for a finite type family over a noetherian integral ring; it
applies because the generic fiber is a curve. Shrink further to make the
geometric fibers integral curves. The generic fiber is geometrically
integral: its base extension to the algebraically closed \(\Omega\) is
integral, so geometric integrality descends. Openness of geometric
integrality for a flat finitely presented family then applies. Only
finitely many equations and this one curve are being spread out.

For each fixed \(D\), the pullback \(\mathcal D\) is invertible and is
flat over \(A\), since \(\mathcal T\) is flat. Both \(\mathcal D\) and
\(\mathcal O_\mathcal T\) are finitely presented. Proper flat
finite-presentation Euler-characteristic constancy (Stacks 0B9T) gives
constant \(\chi(\mathcal D_t)\) and \(\chi(\mathcal O_{\mathcal T_t})\)
on the connected remaining base, and these characteristics commute with
field extension. On an integral proper curve the degree of an invertible
sheaf is their difference. Choose a closed \(k\)-point \(t\) of the
nonempty open base; its residue field is \(k\) by the Nullstellensatz.
The projection formula gives \[\deg a_t^*D=\begin{cases}
e\deg(D|_{a_t(T_t)})&\text{if the image is a curve of generic degree }e,\\
0&\text{if the image is a point}.
\end{cases}\] It is zero when \(D\) is numerically trivial. The same
Euler-characteristic difference computes the degree on \(T\), proving
the assertion. It is not necessary to descend infinitely many \(D\)
simultaneously: the assertion that the numerical kernel is annihilated
is checked one \(D\) at a time. Additivity of degree first gives a
homomorphism on \(\operatorname{Pic}(Z)\), then a \(\mathbb Q\)-linear
map on \(N^1(Z)_\mathbb Q\), and finally a real linear map on the finite
dimensional \(N^1(Z)_\mathbb R\). Every such linear map is continuous.
Thus numerical convergence of classes suffices; Picard convergence is
neither assumed nor needed.

The same argument proves nefness ascends under algebraically closed
field extension: every new curve specializes as above, and a nonnegative
degree remains nonnegative. Nefness descends by extending each old
integral curve and preserving its degree. Degree is unchanged on an
integral curve under field extension, since the cohomology dimensions,
and hence the Euler-characteristic difference, are unchanged. These
arguments also apply to proper curve morphisms, not only embedded
curves. Ampleness of an actual line bundle ascends and descends by fpqc
descent of relative ampleness (Stacks 0D3C); it is not deduced merely
from nefness.

\subsection{3. One graph model for every canonical reference
(M1)}\label{one-graph-model-for-every-canonical-reference-m1}

Fix \(s\in X^{\rm an}\), put \(F=\mathscr H(s)\), and choose a complete
algebraically closed \textbf{nontrivially valued} extension \(K/F\). Let
\(R=K^\circ\) and \(\widetilde K\) be its algebraically closed residue
field. In this section fix a smooth proper connected algebraic curve
\(C/K\) and a morphism \(f:C\to Y_K\) with constant composite to
\(X_K\), equal to the \(K\)-point induced by \(s\). A smooth proper
curve over a field is projective. The maps \(k\to K\) and
\(k\to\widetilde K\) are injective because \(k\) is trivially valued; in
particular both fields have characteristic zero.

Take \textbf{the actual BJ Definition 4.1} representation of a bounded
\(\varphi\in\operatorname{PSH}(Y,\pi^*c_1(L))\):
\[\varphi_i\in\mathcal H_{\mathbb R}^{\rm gf}(L_i),\qquad
\varphi_i\downarrow\varphi,\qquad c_1(L_i)\longrightarrow\pi^*c_1(L).       \tag{3.1}\]
This is a directed net and the \(L_i\in\operatorname{Pic}(Y)_\mathbb Q\)
vary. Every \(\varphi_i\) is finite valued because it dominates the
bounded \(\varphi\). Its finite sections therefore have no common zero,
so \(L_i\) is semiample. This is exactly the bounded clause following
Definition 4.1, not a same-bundle regularization theorem.

Let \(\mathcal Y=Y\times_k R\), the constant projective model. BE
Proposition 5.13 gives the schematic flat graph closure \(\mathcal G\)
of \(f\) in the product of \textbf{one} projective curve model and
\(\mathcal Y\). It is projective and flat over \(R\), dominates that
curve model, and carries a morphism \(\mathcal G\to\mathcal Y\). Models
over a valuation ring are finitely presented by BE Lemma 4.5. Take
relative integral closure in the prescribed generic fiber \(C\). BE
Theorem 4.20 is applicable here because \(K\) is algebraically closed
and nontrivially valued: the closure is finite and its special fiber is
reduced. Its assertion is not used over an arbitrary nondiscrete
valuation ring; BE Remark 4.21 explicitly forbids that generalization.

A semistable model must dominate this chosen graph; existence of an
unrelated semistable model would not suffice. Apply Temkin's
stable-modification theorem (0707.3953v3, Theorems 1.1 and 1.5;
Proposition 4.5.1 and its gluing proof) to the resulting projective
\(R\)-curve with empty marking. Check pure relative dimension one as
follows. The graph and its relative integral closure are integral and
have generic fiber the proper connected smooth \(C\), whose global
functions are \(K\). Stacks 0AY8 applies over the normal integral target
\(\operatorname{Spec}R\): the map is proper, its source reduced, its
generic point dominates, and its generic global functions are \(K\). Its
special fiber is therefore connected. For an actual relatively very
ample bundle, Euler-characteristic constancy (0B9T) applied to each
twist makes its special Hilbert polynomial equal to the generic linear
Hilbert polynomial. Thus its special fiber has maximum dimension one. A
zero-dimensional irreducible component of a reduced finite-type scheme
over an algebraically closed field is an isolated open-and-closed point:
it cannot lie in another component, and there are only finitely many
components. Connectedness excludes such a component. All special
components consequently have dimension one, as do the generic
components. The graph is flat and finitely presented, so it is precisely
a multipointed \(R\)-curve with empty marking in Temkin's definitions in
Section 1.1. The generic fiber is smooth and hence semistable.

For clarity the alteration in Theorem 1.5 does not change our chosen
field. Over \(S=\operatorname{Spec}R\), choose a component of its
generically etale proper alteration \(S'\to S\). Its generic field is
finite separable over \(K\), hence is \(K\) since \(K\) is algebraically
closed. The resulting \(K\)-point extends to a morphism
\(\operatorname{Spec}R\to S'\) by properness. Pulling the stable
modification back along this section gives a semistable \(R\)-curve
\(\mathcal C\) with a proper generic-fiber identity map
\(\mathcal C\to\mathcal G\). It is projective by Temkin 1.1. This is an
actual domination of the chosen graph. Completion at any \(0<|a|<1\)
gives an admissible proper formal model. If a strictly semistable model
is desired, subdivide the finite skeleton by adding type II vertices to
remove loops; BPR 3.13, 3.14 and 4.11 identify these refinements with
dominating formal models. The original morphism to \(\mathcal Y\) is
retained. We use no purported cofinality assertion about an arbitrary
initial semistable model.

All canonical references of the \(L_i\) are now represented on the
\textbf{same} model:
\[\mathcal M_i=\mathcal f^*(L_i\otimes_kR).                              \tag{3.2}\]
The reference on \(L_i\) is the canonical trivial metric over \(k\),
pulled to \(K\). BE 5.10 and 5.13 identify this pullback exactly with
the metric of (3.2). The equality concerns the original bundles and
their trivializations. Since \(L_i\) is semiample, the canonical
reference is also its pure FS metric: a generating system on a divisible
power has maximum norm one in the center frame at each point. This gives
an independent check of both its positivity and its precise reference.
No infinite list of bundles is algebraized on the curve; each already
comes from the one fixed constant \(\mathcal Y\).

Finally the composite \(\mathcal C\to X_R\) factors through the unique
extended section \(\bar s:\operatorname{Spec}R\to X_R\). The generic
constant point extends uniquely by projectivity/separatedness. The
composite and \(\bar s\circ(\mathcal C\to\operatorname{Spec}R)\)
coincide on the generic fiber. The equalizer is closed by separatedness,
and flatness makes that generic fiber schematically dense: on each
affine chart localization into its generic-fiber ring is injective
because the ring is \(R\)-torsion free. Therefore the equalizer is the
whole model as a scheme. This is stronger than topological agreement on
a dense set.

\subsection{4. Fixed background curvature and vanishing coefficients
(M2--M3)}\label{fixed-background-curvature-and-vanishing-coefficients-m2m3}

Let \(E\) run over the finitely many irreducible components of the
reduced special fiber of \(\mathcal C\), and \(\zeta_E\) be their type
II Shilov points. The set \(\Sigma=\{\zeta_E\}\) is fixed independently
of \(i\). Let \(h_{0,i}^C\) denote the actual pulled canonical metric on
\(f^*L_i\). Its raw curvature is
\[\nu_i=c_1(f^*L_i,h_{0,i}^C)=\sum_E a_{i,E}\delta_{\zeta_E},\qquad
a_{i,E}=\deg_{\widetilde K}(\mathcal M_i|_E)\ge0.                       \tag{4.1}\]
This is BE Example 8.11, with its continuous signed-metric extension in
8.10 and the model compatibility of 8.18--8.21; on a curve it agrees
with the CLD curvature current. The degree of a rational bundle is
divided by its cleared tensor denominator. This division must not be
omitted.

Here are the multiplicity and residue conventions in (4.1). The special
fundamental cycle is \([\mathcal C_s]=\sum_E\ell_E[E]\), with \(\ell_E\)
the generic length. A degree over a nonalgebraically closed residue
field already includes the factor \([k_E:\widetilde K]\), where \(k_E\)
is the finite constant field of the proper integral component, relative
to the degree over \(k_E\). Thus an unnormalized cycle formula would
contain \(\ell_E[k_E:\widetilde K]\deg_{k_E}(\mathcal M_i|_E)\). In our
application \(\ell_E=1\) and \(k_E=\widetilde K\): reduced special fiber
and algebraically closed residue field have actually been established.
We do not use \(\operatorname{ord}_E(t)\) over the nondiscrete \(R\),
which has no distinguished uniformizer. In discretely valued
test-configuration formulas the usual multiplicity
\(b_E=\operatorname{ord}_E(t)\) is instead retained. Nilpotents are
never erased when computing a cycle or measure.

The positivity of each coefficient follows also directly:
\(\mathcal M_i|_E\) is the pullback of the semiample
\((L_i)_{\widetilde K}\) through \(E\to Y_{\widetilde K}\) and so has
nonnegative degree. Denote this fixed degree functional on \(N^1(Y)\) by
\(\lambda_E\). Section 2 proves it is well defined and continuous even
though \(E\) and its morphism need not descend to \(k\). The model
factorization in Section 3 says \(\pi\circ f_E\) is constant. Hence
\[\lambda_E(\pi^*c_1(L))=0,\qquad
a_{i,E}=\lambda_E(c_1(L_i))\longrightarrow0.                            \tag{4.2}\]
There are finitely many \(E\), so
\(\nu_i(C^{\rm an})=\sum_Ea_{i,E}\to0\). This is numerical convergence,
not convergence of line bundles. It also agrees with the generic-fiber
degree by flat Euler-characteristic constancy.

Let \(h_i^C=h_{0,i}^C+\varphi_i^C\) be the FS metric obtained by
restricting the actual \(i\)th sections and keeping their finite real
weights. Write \(\mu_i=c_1(f^*L_i,h_i^C)\). It is positive by continuous
FS positivity (BE 7.14, 8.3--8.5). Subtracting the two metrics on
\textbf{the same} \(f^*L_i\) gives the exact distribution identity
\[dd^c\varphi_i^C=\mu_i-\nu_i.                                        \tag{4.3}\]
Its verification is local: in any common frame subtract \(h_{0,i}^C\)
from \(h_i^C\), so transition logarithms cancel. The continuous raw
curve measures represent the corresponding CLD currents; for a
semistable model difference this is Wanner 3.27--3.28, and continuous
dpsh extension follows by the uniform regularizations in BE. The
backgrounds do not silently switch to a fixed ample metric. Both
\(\mu_i\) and \(\nu_i\) have the same raw degree, but only \(\nu_i\)
needs the fixed support.

\subsection{5. Correct Green sign and passage of the decreasing
net}\label{correct-green-sign-and-passage-of-the-decreasing-net}

Take a connected strict affinoid domain \(V\) in \(C^{\rm an}\) with
nonempty boundary disjoint from \(\Sigma\). For each
\(\zeta\in\Sigma\cap V\), Thuillier 3.3.7(ii) provides a continuous
lisse Green function \(g_\zeta^V\), zero on \(\partial V\), with
\[dd^cg_\zeta^V=-\delta_\zeta\quad\text{on }V\setminus\partial V.
\] The boundary measure is not discarded globally. Put
\[p_i=-\sum_{\zeta\in\Sigma\cap V}a_{i,\zeta}g_\zeta^V.                 \tag{5.1}\]
Then \(dd^cp_i=\nu_i\) on the interior and
\[\|p_i\|_V\le\sum_{\zeta\in\Sigma\cap V}a_{i,\zeta}\|g_\zeta^V\|_V
\longrightarrow0.                                                    \tag{5.2}\]
This proves uniform zero correction. It uses \textbf{fixed}, bounded
Green functions and a finite support. A small mass with moving poles
would not prove (5.2).

The functions \(u_i=\varphi_i^C+p_i\) are continuous and have positive
CLD current on the interior by (4.3). All hypotheses of Wanner Theorem
4.7 hold: \(K\) is algebraically closed, complete and nontrivially
valued; \(C\) is a smooth proper algebraic curve; the tested set is an
ordinary open; and the functions are finite continuous there. They are
therefore Thuillier subharmonic. They need not decrease.

We spell out the net argument. Let \(q\) be continuous and harmonic on
\(V\) and suppose \(\varphi^C\le q\) on \(\partial V\). For
\(\epsilon>0\), the closed boundary sets
\[\{x\in\partial V:\varphi_i^C(x)\ge q(x)+\epsilon\}\] decrease and have
empty intersection. Compactness and directedness give one index after
which all are empty. At the same time (5.2) gives \(|p_i|<\epsilon\).
Subharmonic comparison yields \(u_i\le q+2\epsilon\) on \(V\), and hence
\(\varphi_i^C\le q+3\epsilon\). Take the decreasing original limit and
then \(\epsilon\downarrow0\). This is finite-intersection boundary
comparison, not Dini for a possibly discontinuous limit and not monotone
convergence of the corrected functions.

There is a simple way to cover also test domains whose boundaries meet
\(\Sigma\), without asserting an unproved restricted-basis criterion.
For any proper strict affinoid \(W\), choose a type II point \(\eta\)
outside \(W\). Thuillier 3.3.7(i) gives bounded global lisse functions
\(g_{\zeta,\eta}\) with \(dd^cg_{\zeta,\eta}=\delta_\eta-\delta_\zeta\).
Use \(p_i'=-\sum_{\zeta\in\Sigma}a_{i,\zeta}g_{\zeta,\eta}\). Its norm
goes to zero on the proper curve and its curvature is
\(\nu_i-(\sum a_{i,\zeta})\delta_\eta\). Thus \(\varphi_i^C+p_i'\) is
subharmonic on a neighborhood of \(W\). The identical boundary argument
proves harmonic comparison on every such \(W\). These are exactly
Thuillier's comparison tests. The bounded usc function \(\varphi^C\) is
consequently subharmonic on the whole curve and is not identically
\(-\infty\).

On a proper connected curve it attains a finite maximum. Thuillier
3.1.11 gives local constancy at any local maximum. The maximum set is
closed by usc and open by that principle, hence is the whole connected
curve. Therefore
\[\varphi\circ f^{\rm an}\quad\text{is constant on }C^{\rm an}.          \tag{5.3}\]
This establishes the curve-constancy consequence of M1--M3 on the actual
Definition 4.1 net. No affine regularization over \(K\), no same-class
FS net, and no Picard convergence has been used.

\subsection{6. All analytic fiber points have the same
value}\label{all-analytic-fiber-points-have-the-same-value}

First justify geometric connectedness. The finite normalization
\(\nu:X^\nu\to X\) is a universal homeomorphism by Stacks 0GIQ, since
\(X\) is irreducible geometrically unibranch. The normal-target map
\(\rho:Y\to X^\nu\) satisfies all six assumptions of Stacks 0AY8: it is
proper; its target is integral and normal; its source is reduced; its
single generic point maps to the target generic point; and its generic
fiber has \(H^0=k(X)\). The normal-target Stein factor is therefore the
target, and every fiber is geometrically connected.

For any field point \(\operatorname{Spec}F\to X\) the normalization
fiber is finite radicial. In characteristic zero its unique residue
field is \(F\), so its reduction is \(\operatorname{Spec}F\), though it
may have nilpotents. Reducing after an algebraic closure gives the same
underlying topological fiber of \(\rho\). Hence every geometric fiber of
\(\pi\) is connected. This reduction is used only for connectivity and
curves. It is never used to replace a fundamental cycle in a measure
identity.

Fix \(s\in X^{\rm an}\) and \textbf{one} closed point \(z\) of the
nonempty projective \(F\)-fiber \(Y_s\), where \(F=\mathscr H(s)\). Its
residue field \(F_z/F\) is finite. The complete norm of \(F\) extends
uniquely to \(F_z\), giving a fixed analytic point \(z^{\rm an}\) over
\(s\). Now let \(y\) be any analytic point over \(s\). Embed
\(\mathscr H(y)\) isometrically into a complete algebraically closed
nontrivially valued field \(K_y\). If its norm was trivial, first adjoin
a Gauss variable with radius \(0<r<1\) and complete an algebraic
closure. Choose an \(F\)-embedding of \(F_z\) into \(K_y\); the unique
norm-extension property makes this embedding isometric. This is done
separately for each \(y\), with the same original \(z\). It provides
\(K_y\)-rational lifts of \(y\) and \(z^{\rm an}\) to the geometric
fiber. We do not descend a family of analytic points to a common field.

In a connected reduced projective fiber over algebraically closed
\(K_y\), the finite graph of irreducible components joined by nonempty
intersections is connected. Choose closed points in each intersection
along a path between components containing the two lifts. Within one
component \(T\), use a projective resolution \(\widetilde T\to T\) in
characteristic zero, and choose points above the two tested points. Here
is an explicit smooth curve-morphism connection. If \(\dim T=0\) there
is no step; if it is one the smooth proper resolution itself suffices.
If its dimension is \(d\ge2\), blow up the two distinct chosen points on
the smooth projective \(\widetilde T\). Choose a very ample bundle on
this smooth connected blowup. Iteratively take \(d-1\) general
hyperplane sections. Stacks 0FD6 gives smoothness (generated bundle and
an immersion); 0FD9 gives connectedness at each stage, since the
intermediate smooth scheme has dimension at least two and hence depth at
closed points at least two. Smooth connected schemes here are
irreducible. Impose also proper intersection with each exceptional
divisor, an open condition on the finite tuple of hyperplanes. The final
smooth connected proper curve meets each exceptional divisor: its
intersection number with it is the positive top self-intersection of the
ample restriction on that \(\mathbb P^{d-1}\). Its composite map to
\(T\) therefore passes through both original points. This is a morphism,
not an assumption that an embedded curve in the singular component is
smooth.

The resulting finite chain consists of smooth proper connected curve
morphisms into \(Y_{K_y}\) whose composites to \(X_{K_y}\) are constant.
Section 5 applies to each. It follows that
\(\varphi(y)=\varphi(z^{\rm an})\) for every \(y\) in the fiber, for
bounded genuine \(\varphi\).

For arbitrary genuine \(\varphi\in\operatorname{PSH}(Y,\pi^*L)\), the
floors \(\max(\varphi,-m)\) remain genuine in the \textbf{same nef
class}: constants are psh in a nef class by BJ 4.3, and finite maxima
are stable by BJ 4.7. Apply the bounded result and let \(m\to\infty\).
This proves actual equality of values, including \(-\infty\), on each
analytic fiber. Density plus usc is not used to determine those values.

Define \(\psi(s)\) to be this common value. It is usc: for every real
\(a\), the superlevel \(\{\psi\ge a\}\) is the image under the proper
analytic map \(\pi^{\rm an}\) of the closed superlevel
\(\{\varphi\ge a\}\). Proper maps between these compact Hausdorff spaces
are closed. It is generically finite since the generic trivial valuation
upstairs maps to the generic trivial valuation downstairs. We have
proved scalar descent and usc; genuine original-class membership is the
next gate, not a consequence of these topological assertions.

\subsection{7. Exact original reference requirement and
C1}\label{exact-original-reference-requirement-and-c1}

For a frame \(e\) of the original \(L\) on \(U\subset X\), write
\(h_{0,e}=-\log\|e\|_0\) for its canonical reference weight. The
descended metric must have the weight \[b_e=\psi+h_{0,e},\qquad
b_e=b_{e'}-\log|e/e'|,\qquad
\pi^*b_e=\varphi+\pi^*h_{0,e}.                                        \tag{7.1}\]
The pulled reference is exactly the canonical reference on \(\pi^*L\),
by pullback of the constant model; it is not one of the varying \(L_i\)
unless they actually equal \(\pi^*L\). Formula (7.1) proves cocycle
compatibility, but Definition 4.1 alone supplies no convergence of the
different \(h_{0,i}\) to \(\pi^*h_{0,e}\). We now prove the actual local
contact application by changing the approximating classes in a
controlled, justified way.

\textbf{Original-reference contact lemma.} Fix a very ample actual
bundle \(A\) on \(Y\), with its canonical metric. Write
\(\theta=\pi^*c_1(L)\). The function \(\varphi\) is already genuine on
\(Y\), so actual BJ 12.18(ii) applies: it gives a decreasing sequence
\(a_j\in H(M_j)\), with \(M_j\) ample rational bundles,
\(c_1(M_j)\to\theta\), and \(a_j\downarrow\varphi\). This source theorem
has varying bundles and does not assume EP on \(X\). Its countable
regularization proof uses the already constructed genuine energy theory;
it is not a first-membership theorem and is not regularization on the
nef \(\pi^*L\).

Put \(\delta_r=2^{-r}\). By numerical convergence and openness of the
ample cone, choose an increasing subsequence \(j(r)\) so that
\[\delta_r c_1(A)-(c_1(M_{j(r)})-\theta)\text{ is ample}.                 \tag{7.2}\]
Indeed, the open neighborhood \(\delta_r c_1(A)-\operatorname{Amp}(Y)\)
of zero contains the numerical difference eventually. Constants are psh
in this nef difference, and BJ 4.7 therefore gives
\(a_{j(r)}\in\operatorname{PSH}(Y,B_r)\) for the \textbf{actual} ample
rational bundle \[B_r=\pi^*L+\delta_r A.\] This class change is proved
by adding the constant zero; there is no identification of \(M_{j(r)}\)
with \(B_r\) or claim about their Picard distance. Each \(a_{j(r)}\) is
finite continuous. BJ 4.16 and Example 3.9 now give \(q_r\in H(B_r)\)
with rational weights and \(\|q_r-a_{j(r)}\|_\infty<\epsilon_r=4^{-r}\).
These are approximations in the same actual \(B_r\) at the fixed \(r\),
whose existence follows from continuous genuine membership in that ample
class. The bundles \(B_r\) vary.

Set \(c_r=2\sum_{l\ge r}\epsilon_l=\frac83\,4^{-r}\). The difference
\[q_r+c_r-q_{r+1}-c_{r+1}\ge2\epsilon_r-\epsilon_r-\epsilon_{r+1}>0\]
shows \(q_r+c_r\downarrow\varphi\). The limit is unchanged because the
uniform approximation errors and shifts go to zero. All shifts are
rational.

Cover \(Y\) by the finite section frames \(s\) of a generating system of
the very ample \(A\). On \(D(s)\) its canonical reference weight is
\(h_{A,s}=\max_l\log|s_l/s|\ge0\), with the branch \(s/s=1\). On
\(\pi^{-1}U\cap D(s)\), tensor multiplicativity of the actual canonical
metrics gives the reference of \(B_r\) exactly as
\(\pi^*h_{0,e}+\delta_r h_{A,s}\). Consequently the weights
\[w_r=q_r+c_r+\pi^*h_{0,e}+\delta_r h_{A,s}                            \tag{7.3}\]
decrease to \(\varphi+\pi^*h_{0,e}\) on this one common ordinary open:
both variable terms decrease and \(h_{A,s}\ge0\) is finite there. Each
\(w_r\) is precisely the local weight of its rational FS metric on
\(B_r\), plus the rational constant \(c_r\). Clear its tensor
denominator; in the frame \(e\otimes s^{\delta_r}\) it is a finite
maximum of rational multiples of logarithms of regular section
coefficients, plus rational constants, in the \textbf{original ring of
this witness chart}. The generating sections have no common zero. At any
tested point choose a maximizing nonzero branch. Its rational affine
expression lies below the maximum and equals it with the same constant.
Branches vanishing there are inactive and can be removed on a smaller
point neighborhood; a sufficiently low rational floor is inactive on a
compact smaller neighborhood. Thus the supplied exact rational contact
condition holds, including its finite-floor convention. The functions
for the entire sequence are defined on the common chart; unit reductions
may be made separately when testing each continuous approximant. For any
finite tuple the intersections give one common ordinary neighborhood.

This proves that every original-frame resolved weight is locally a
decreasing rational contact limit on the \textbf{one} witness \(Y\). It
proves the required R-original application, including unbounded
\(\varphi\), since BJ 12.18(ii) gives finite continuous approximants
even then. Alternatively bounded floors stay in the same nef class by
Section 6. No field extension or affine regularization over a completed
residue field occurs in this construction. The M1--M3 proof still used
the actual Definition 4.1 net; this separate contact lemma does not
replace it by a fictitious same-bundle net.

The sole remaining analytic hypothesis is now stated in weak-to-genuine
form:

\textbf{C1-original (conditional).} Over every countable algebraically
closed trivially valued characteristic-zero field, an original
geometrically unibranch projective \(X\) with actual ample rational
\(L\) has the following global membership theorem. If a generically
finite usc scalar \(\psi\) has original weights \(\psi+h_{0,e}\) with
the actual \(L\) cocycle, and their pullbacks on one proper birational
smooth witness are locally decreasing rational contact limits in the
exact task-granted sense, then \(\psi\in\operatorname{PSH}(X,L)\) in BJ
Definition 4.1. All required geometric
compactification/domination/resolution-witness antecedents must be
included in that theorem's scope. No completed-field version is required
or asserted.

Sections 6 and (7.1)--(7.3) verify these admissibility and reference
hypotheses for our descended scalar. If the countable theorem also
formulates an upper-growth condition, it is present here:
\(\sup\psi=\sup\varphi<\infty\) by the generic-trivial-value formula for
genuine potentials, and in the section frame of a sufficiently divisible
very ample original \(L\) the reference \(h_{0,e}\) is a finite affine
FS maximum in the original section ratios. Thus
\(b_e\le h_{0,e}+\sup\psi\) has the actual original upper-growth
reference, with generators appended at dominated weights as in the
task-granted localization construction. One may first apply C1 to the
bounded scalar floors and then use genuine decreasing closure for the
unbounded limit. C1 supplies the remaining global membership conclusion,
not a smooth or changed-class conclusion. This hypothesis is neither
proved here nor inferred from Wu's printed result or independent Runge
work.

Assuming C1-original, pullback is onto. Its forward inclusion is BJ 4.7;
injectivity is analytic surjectivity. The divisorial sets on \(X\) and
\(Y\) are the identical valuations of their common function field
(including the trivial valuation), not separate dense sets. Their BJ
topologies are pointwise convergence on these sets; hence the bijection
is a homeomorphism. Analytic surjectivity also gives
\(\sup_X\psi=\sup_Y\pi^*\psi\). Consequently normalized PSH spaces are
homeomorphic. Smooth nef BJ 5.20--5.21 gives compactness of normalized
\(\operatorname{PSH}(Y,\pi^*L)\), even though \(\pi^*L\) need not be
ample. Target normalized compactness follows, and BJ 5.11 gives EP on
this \textbf{countable original target}. BJ 5.13 has not been used here.
Using it before approaching target EP would be circular.

Thus the explicitly stated original-class weak-to-genuine C1 implies
countable geometrically unibranch ample EP. The local-reference
application has been proved above; C1 itself remains conditional.

\subsection{8. Finite data and arbitrary
cardinality}\label{finite-data-and-arbitrary-cardinality}

Assume henceforth countable EP just obtained conditionally. We
reconstruct the full finite-data route. A finite field of definition
means a finitely generated subfield of \(k\) containing the finitely
many coefficients in the chosen presentations, followed by its algebraic
closure inside \(k\). This is countable in characteristic zero. It need
not contain the real constants occurring as weights; those are kept as
the same finitely many real numbers, external to the coefficient field.

For finite algebraic data, descent is literal. Projective equations,
finite affine presentations, transition functions of actual line
bundles, finite section systems and their equalities, flag ideals, graph
maps, the \(\mathbb G_m\)-action and its identities, and the specified
generic-fiber isomorphism all involve finitely many coefficients and
identities. Enlarge the countable subfield to contain them. Generation
by the descended sections is detected by the coherent cokernel, which is
zero after faithfully flat extension iff it was zero. An
integral/reduced projective scheme downstairs is detected in the same
manner; since the field is algebraically closed and characteristic zero
it is geometrically integral/reduced. Actual ample bundles and relative
ample model bundles descend by fpqc ampleness. A finite birational map
is descended together with a dense open on which it is an isomorphism;
finiteness and that isomorphism are fpqc properties. The dense open
remains nonempty since its extension is nonempty.

Descend the normalization as a \textbf{finite map with its actual
source}, not as an unsupported slogan. Include its finite presentation,
algebra inclusion and dense-open isomorphism. Its source is normal:
normality descends through the faithfully flat field extension. Over an
algebraically closed characteristic-zero field it is geometrically
normal (Stacks 0BXT, since the field is already perfect). The finite
birational normal source is therefore the normalization and remains so
after any field enlargement. Its universal-homeomorphism property also
descends: after extending algebraically closed fields, a finite map's
geometric fibers and residue extensions detect universal injectivity,
while surjectivity is fpqc. Equivalently the finite diagonal's
surjectivity can be checked after this base change. Hence the original
downstairs \(X_0\) is geometrically unibranch; one has not silently
replaced a nonnormal generic fiber by its normalization.

For a finite model obstacle or Hdom competitor descend its
\textbf{actual generic-fiber} model and bundle. Flatness over the test
line is detected faithfully flat, or directly by torsion freeness over
\(k_0[t]\). For relative integral closure use BJ 1.13 and 1.15: it is
equivalent to vertical \(R_1\) and to the special fiber having no
embedded associated points (vertical \(S_2\)). Algebraically closed
characteristic-zero field extensions are regular. In local noetherian
rings the flat depth formula and the corresponding dimension formula,
with regular closed fibers, preserve and reflect \(S_1\); regularity
preserves and reflects the regular local rings at the generic points of
the central fiber. Indeed all central components over an algebraically
closed field are geometrically integral, so their generic points remain
precisely generic central points. BJ's two vertical conditions therefore
descend and ascend. This uses only rings finite type over fields, not a
false noetherian assertion for \(K^\circ\).

Alternatively take relative integral closure of a descended raw model
\textbf{inside its prescribed original generic fiber} and apply those
same two vertical conditions to see that its base extension is the
prescribed relatively integrally closed model. The closure is finite
because these models are finite type over a field. Absolute
normalization would change a nonnormal generic fiber and is not allowed.

Central cycle data are retained. At each generic central point the local
ring is a DVR, and \(b_E=\operatorname{ord}_E(t)\) equals the generic
length in the special fiber. Flat field extension preserves this length
(tensor a composition series, whose residue-field factors remain fields
at the corresponding generic point). Each component is geometrically
integral; intersection degrees of its model bundles are unchanged, by
flat cohomological base change for multivariate Hilbert polynomials.
Consequently relative nefness is preserved and reflected using the
curve-specialization lemma of Section 2 on each projective special
component, and relative ampleness uses fpqc descent. The finite divisor
function, finite flag-ideal formula and finite FS section maximum are
the \textbf{same scalar functions after analytic base extension}:
valuation evaluation and the pulled constant references commute. No
algebraic field is asked to encode a real weight.

These statements suffice for rational PL obstacles and finite Hdom
competitors on an actual rational ample \(L\). For a finite real PL
obstacle one can either retain its finite real coefficients or
approximate it uniformly by rational PL obstacles; only the rational
case is needed for EP. An arbitrary real numerical class is handled
later and is not incorrectly packaged as one actual rational bundle.

\subsection{9. Raw continuous MA under complete field
extension}\label{raw-continuous-ma-under-complete-field-extension}

Let \(k_0\subset k\) be algebraically closed trivially valued fields,
and \(p:X_k^{\rm an}\to X_0^{\rm an}\) the analytic projection. It is
continuous and surjective. Surjectivity follows by extending the valued
residue field of any analytic point to a complete valued field
containing \(k\); equivalently it is analytic proper base-change
surjectivity. For finite continuous genuine metrics on actual ample
rational bundles, \[p_*\left[\bigwedge_{j=1}^n
 (c_1(L_{j,k})+dd^cp^*\phi_j)\right]
=\bigwedge_{j=1}^n(c_1(L_j)+dd^c\phi_j),                               \tag{9.1}\]
and the raw mass on either side is \((L_1\cdots L_n)\).

Here is the actual operator calibration. BE Proposition 8.3(iv) states
complete-field invariance for smooth metrics \textbf{with the analytic
fundamental cycle}, Theorem 8.4 uniquely extends under uniform
convergence, and 8.5 retains positivity and raw intersection mass.
Semiample continuous FS-limit metrics are psh-regularizable by 7.14. BJ
4.16 supplies uniform FS approximation for continuous genuine potentials
on an actual ample rational bundle, and BJ 7.23 identifies its raw
measure with this continuous theory on the general reduced projective
variety. Therefore (9.1) follows by uniform approximation and
pushforward against every continuous test. This argument applies on the
original singular or nonnormal \(X_0\), not only its normalization.

The raw model formula provides a second calibration:
\[\sum_E b_E\big((\mathcal L_1|_E)\cdots(\mathcal L_n|_E)\big)\delta_{v_E}.
\] The cycle multiplicities and degrees are those retained in Section 8;
\(v_E=b_E^{-1}\operatorname{ord}_E|_{k(X)}\). For possibly nonreduced
schemes, the analytic fundamental cycle is instead the sum of generic
lengths times the component cycles, and finite residue degrees enter
each degree. Formula (9.1) preserves those factors, as BE 8.1--8.5
explicitly records. We do not apply it to an arbitrary discontinuous
metric by a formal weak-limit definition. Only continuous metrics are
used in the envelope descent; the genuine unbounded products are
supplied separately by BJ 7.18--7.19.

\subsection{10. Orthogonality bounds every finite
competitor}\label{orthogonality-bounds-every-finite-competitor}

Fix a rational \(f\in\operatorname{PL}(X)\) and the actual ample
rational \(L\). Descend these and normalization to a countable
algebraically closed \(k_0\) as in Section 8. Write \(f=p_0^*f_0\).
Countable EP gives the finite continuous genuine envelope
\(\phi_0=P_{L_0}(f_0)\); a constant below \(\min f_0\) is a competitor,
so it is not \(-\infty\). Put \(\phi=p_0^*\phi_0\). Base extension of
the Definition 4.1 net gives actual genuine membership of this
continuous metric on \(L\); the numerical pullback map is linear
continuous by Section 2. The envelope satisfies \(\phi\le f\).

Countable EP alone gives the two tools needed here, without importing
the desired inverse equation. Orthogonality follows directly from BJ
13.1 under its explicit EP premise. For a continuous obstacle, its short
proof uses the genuine energy derivative: if \(g=f_0-\phi_0\ge0\), then
\(P(\phi_0+tg)=\phi_0\) for \(0\le t\le1\), so the derivative of
\(E\circ P\) is zero and \(\int g\,\mathrm{MA}(\phi_0)=0\). Domination
is genuine BJ 10.6 and needs no EP or inverse Dirac equation. Its input
is a pointwise a.e. inequality for the genuine MA of one finite-energy
potential.

Take any finite Hdom competitor \(h\) with \(h\le f\). Its finite model
data may require a \textbf{new} countable algebraically closed
\(k_1\supset k_0\). Descend them there and write \(h=p_1^*h_1\).
Surjectivity of \(p_1\) gives \(h_1\le f_1\). The pullback \(\phi_1\) of
\(\phi_0\) is genuine continuous. By (9.1), the nonnegative continuous
integrand gives
\[\int_{X_1^{\rm an}}(f_1-\phi_1)\,\mathrm{MA}^{\rm raw}(\phi_1)
=\int_{X_0^{\rm an}}(f_0-\phi_0)\,\mathrm{MA}^{\rm raw}(\phi_0)=0.         \tag{10.1}\]
Thus \(f_1=\phi_1\) a.e. for that measure, and \(h_1\le\phi_1\) a.e.
Domination gives \(h_1\le\phi_1\) everywhere. Pull back to \(k\) to get
\(h\le\phi\). One can also transfer (10.1) directly to \(k\); the
separate \(k_1\) makes explicit that each finite competitor has its own
field. No equality between a support and an inverse image of a support
is used, and no entire family is descended.

Finally let \(q\in\operatorname{PSH}(X,L)\) be \textbf{any} competitor
with \(q\le f\). BJ 4.15(iii), not an unsupported same-bundle
regularization assertion, gives a decreasing net
\(h_i\in\mathrm{Hdom}(L)\) with \(h_i\downarrow q\). For rational
\(\epsilon>0\), the closed sets \(\{h_i\ge f+\epsilon\}\) on compact
\(X^{\rm an}\) decrease with empty intersection. Directed finite
intersection gives eventual emptiness. Then \(h_i-\epsilon\le f\); this
remains a finite Hdom function (rational constant shift corresponds to a
rational central divisor). The preceding bound applies separately to
each such function, giving \(h_i-\epsilon\le\phi\). Take its decreasing
limit and then \(\epsilon\downarrow0\) to obtain \(q\le\phi\). Therefore
the exact, unregularized envelope equality is
\[P_L(f)=p_0^*P_{L_0}(f_0).                                             \tag{10.2}\]
It is continuous and genuine. The proof bound every arbitrary competitor
through finite feasible majorants; it did not suppose an arbitrary psh
net had a common countable field of definition.

PL density and the elementary inequality
\(\|P(f)-P(g)\|_\infty\le\|f-g\|_\infty\) extend continuity and genuine
membership to every finite continuous obstacle: approximate uniformly by
rational PL, and use uniform closure of PSH. BJ 5.17 gives EP for the
rational ample class. For every real ample numerical class \(\alpha\),
choose rational ample \(\omega_j\to\alpha\) with \(\omega_j-\alpha\)
ample, hence nef. Such a sequence exists in the finite-dimensional
numerical space by openness of the ample cone and a small positive
displacement toward a fixed rational ample direction. Constants are psh
in each nef difference, so
\(\operatorname{PSH}(\alpha)\subset\operatorname{PSH}(\omega_j)\) by BJ
4.7. For any bounded-above family of \(\alpha\)-psh functions, its usc
supremum is \(\omega_j\)-psh by EP there. BJ varying-class closure makes
it \(\alpha\)-psh. This proves EP for all ample numerical classes over
arbitrary-cardinality \(k\), conditional on countable EP, with no
missing operator premise in the continuous transfer (9.1).

\subsection{11. Normalization is onto only after target
EP}\label{normalization-is-onto-only-after-target-ep}

The finite birational universal homeomorphism \(\nu:X^\nu\to X\) induces
a homeomorphism on analytifications: each finite radicial fiber has just
one analytic point, since its unique residue-field norm is fixed, and a
proper continuous bijection of compact Hausdorff spaces is a
homeomorphism. This is a topological assertion and does not alone give
genuine psh descent.

There is an explicit conductor blowup realizing \(\nu\). Write
\(B=\nu_*\mathcal O_{X^\nu}\) and
\(J=\operatorname{Ann}_{\mathcal O_X}(B/\mathcal O_X)\). It is a nonzero
coherent ideal: clearing the denominators of finitely many integral
generators of the finite birational algebra produces a nonzero conductor
element. Choose an ample \(A\) and a nonzero section
\(a\in H^0(X,J\otimes A)\) after sufficiently large twisting. In a frame
of \(A\) let \(a_U\) be its coefficient and put
\(I|_U=a_UB|_U\subset\mathcal O_U\). These ideals agree on overlaps
because the frame ratios are units of \(\mathcal O_X\), hence of \(B\).
They form a coherent nonzero ideal of \(\mathcal O_X\).

On an affine integral chart \(A_0\subset B_0\), \(I=aB_0\) and
\(I^m=a^mB_0\) for \(m\ge1\). Its Rees algebra is
\(A_0\oplus\bigoplus_{m\ge1}a^mB_0t^m\). The single chart \(D_+(at)\)
covers its Proj. To see this, every degree-one generator \(abt\) with
\(b\) integral over \(A_0\) has a monic equation whose coefficients in
degree \(j\) contain \((at)^j\); a homogeneous prime containing \(at\)
consequently contains all degree-one generators, hence the irrelevant
ideal, and is not a Proj point. In the localization the degree-zero ring
is exactly \(B_0\), since \((abt)/(at)=b\) and all degree-zero fractions
are in \(B_0\). Thus \(\operatorname{Bl}_I(X)=X^\nu\)
scheme-theoretically, and these chart identities glue. This calculation
does not use a conductor ideal with incorrectly asserted invertible
pullback.

Now, \textbf{after} Section 10 proves target EP conditionally for every
ample numerical class, choose target
\(\theta_m=(1+m^{-1})c_1(L)>c_1(L)\). BJ 5.13 applies and gives the
genuine homeomorphism \[\nu^*:\operatorname{PSH}(X,L)\xrightarrow{\sim}
\operatorname{PSH}(X^\nu,\nu^*L).                                      \tag{11.1}\]
The conductor calculation also justifies its blowup presentation in the
lemma's proof. This application is not used to obtain the first target
EP. Actual suprema agree by surjectivity, and continuous functions
correspond because \(\nu^{\rm an}\) is a homeomorphism.

\subsection{12. Exact mixed products, energy, and finite-energy
transport}\label{exact-mixed-products-energy-and-finite-energy-transport}

The general-variety scope of BJ is essential. Its initial conventions
allow reduced nonnormal varieties, and Section 7 resumes that scope. The
genuine pairing and operator already exist there, independently of EP.
Let \(V=(L^n)>0\) and use the raw operator
\[\mathrm{MA}^{\rm raw}_L(\psi)=(c_1(L)+dd^c\psi)^n,
\qquad \mathrm{MA}_L=V^{-1}\mathrm{MA}^{\rm raw}_L.                        \tag{12.1}\]
The normalization is a calibration, not a definition of a different
inverse measure.

BJ Proposition 7.2, equation (7.5), gives equality of the genuine mixed
energy pairings under normalization pullback; Proposition 7.19, equation
(7.18), gives \[\nu_*\bigwedge_j(c_1(\nu^*L_j)+dd^c\nu^*\psi_j)
=\bigwedge_j(c_1(L_j)+dd^c\psi_j)                                      \tag{12.2}\]
for genuine signed finite-energy entries in its stated scope. This
already includes unbounded entries; it is not inferred by taking an
arbitrary formal weak limit. At the model level normalization has degree
one, sends its fundamental cycle to \([X]\), and is an isomorphism at
generic vertical DVRs, so \(b_E\) and the projected intersection
coefficients agree. For general entries the source extends the actual
symmetric pairing using monotone energy continuity and defines the
signed Radon measure by testing rational PL; density identifies the
measure. Both normalizations in (12.1) use the same \(V\), by the
degree-one projection formula.

Put \[E_L(\psi)=\frac{\langle(c_1(L),\psi)^{n+1}\rangle}{(n+1)V},\qquad
\mathcal E^1(X,L)=\{\psi\in\operatorname{PSH}(X,L):E_L(\psi)>-\infty\}.
\] For arbitrary genuine psh inputs the energy is the decreasing
extension of the bounded energy. Floors pull back exactly, and their raw
pairings are equal by 7.2. Monotone continuity therefore gives
\[E_L(\psi)=E_{\nu^*L}(\nu^*\psi),                                      \tag{12.3}\]
including the value \(-\infty\). Once (11.1) is known, this proves onto
\(\mathcal E^1\) and its weak and strong topological compatibility: the
weak topologies use the identical divisorial valuations, and the strong
ones add continuity of this same energy. There is no use of I-Cauchy
completion to manufacture an operator or a potential.

The target measure is a positive Radon measure \(\eta\) with \textbf{raw
mass \(V\)}, satisfying
\[E_L^\vee(\eta/V)=\sup_{\psi\in\mathcal E^1(X,L)}
\left(E_L(\psi)-\int\psi\,d(\eta/V)\right)<\infty.                       \tag{12.4}\]
BJ 9.20 on irreducible general \(X\) equates this with integrability of
every genuine \(\mathcal E^1\) potential. No skeleton or support
restriction appears. Transport \(\eta\) through the analytic
normalization homeomorphism to \(\eta^\nu\). Equations (11.1)--(12.3)
and the ordinary change-of-variables formula give equal dual energies,
equivalent integrability conditions and exact mass. Formula (12.2) then
transports actual finite-energy solutions and all mixed products.
Nilpotent cycle factors from a model or a different field are retained
throughout; the only topological reductions were in Section 6.

\subsection{13. Conditional variational equation, with all
replacements}\label{conditional-variational-equation-with-all-replacements}

Conditional on EP obtained above, the actual general-variety theorem is
BJ \textbf{12.8}, whose equivalence includes onto finite-energy MA. We
record the proof ingredients rather than substituting the smooth
corollary 12.9.

\begin{itemize}
\tightlist
\item
  Compactness is target normalized compactness, equivalent to EP by BJ
  5.11. On countable targets it was transferred from smooth \textbf{nef}
  compactness only under C1-original. On arbitrary-cardinality targets
  Section 10 proves EP by finite obstacle and finite competitor descent.
  No metrizable compactness or countable net was silently assumed.
\item
  Envelopes are genuine on the original variety and original ample
  bundle by (10.2), PL density, varying-class closure and BJ 5.17.
  Orthogonality uses the nonnegative integral (10.1). The new Wu v1
  envelope proof is not an independent proof of this gate: its
  LF2697--2755 explicitly invokes resolved-nef descent and its
  weak/global theorem.
\item
  Energy differentiation is BJ 8.5--8.6 for the \textbf{outer}
  decreasing energy \(E^\downarrow\), on the reduced-projective general
  scope. For PL \(g\) it gives
  \(E^\downarrow(\psi+tg)=E(\psi)+t\int g\,\mathrm{MA}_L(\psi)+O(t^2)\);
  uniform PL density and the Lipschitz bound give the continuous-test
  derivative. Under EP, \(E^\downarrow(\psi+tg)=E_L(P_L(\psi+tg))\). The
  derivative concerns genuine energy and the genuine operator, not an
  auxiliary law.
\item
  Actual attainment uses a normalized maximizing sequence for (12.4). It
  has \(J_{\eta/V}\to0\). BJ's bound (9.11) makes \(J\) eventually
  bounded. Weak normalized compactness gives a subnet; lower
  semicontinuity of \(J\) keeps its limit in \(\mathcal E^1\). BJ 9.19
  gives convergence in \(L^1(\eta/V)\) on this bounded-energy set;
  energy usc then gives an actual maximizer. Without compactness the
  same maximizing sequence would not prove existence.
\item
  To identify its equation let \(\mu=\eta/V\), \(\psi\) be that
  maximizer and \(g\in C(X^{\rm an})\). The envelope is squeezed between
  \(\psi-|t|\|g\|_\infty\) and \(\psi+|t|\|g\|_\infty\), so is genuine
  finite energy under EP. For both signs of \(t\),
  \(P_L(\psi+tg)\le\psi+tg\), so maximality gives
  \(E_L(P_L(\psi+tg))-\int(\psi+tg)d\mu\le E_L(\psi)-\int\psi d\mu\).
  The outer derivative yields
  \(\int g\,\mathrm{MA}_L(\psi)=\int g\,d\mu\) for every continuous
  test. Thus \(\mathrm{MA}^{\rm raw}_L(\psi)=\eta\). This is equality of
  actual Radon measures.
\item
  Finite-energy comparison is genuine BJ 7.40--7.44 on the general
  variety, obtained from winning-locus locality, maxima, equal raw mass
  and decreasing continuity. BJ 10.4 gives uniqueness modulo one real
  constant for equal MA; its \(I=0\) consequence uses genuine
  finite-energy estimates and divisorial order, not density plus usc. BJ
  10.6 supplies the domination already used in Section 10.
\end{itemize}

The raw objective \(VE_L(\psi)-\int\psi d\eta\) changes under addition
of \(c\) by \(c(V-\eta(X^{\rm an}))\), so the mass in (12.4) is
necessary for invariance. At positive ample volume the probability
calibration (12.1) is valid. Nothing in this argument asserts the
zero-volume or nonample axes.

Consequently: \textbf{if C1-original holds}, the foregoing source-based
applications give EP and the genuine finite-energy equation for every
variety in the retained arbitrary-cardinality domain, with exact mass
(12.4), arbitrary support, and uniqueness modulo constants. The
reference/contact application was proved in Section 7 using varying
actual ample bundles. The manuscript does not prove C1-original, so this
conditional statement does not satisfy the unchanged original completion
criterion. No independent Runge result is assumed, and candidate=false.

\subsection{14. Documentary observations, failures, and
rerun}\label{documentary-observations-failures-and-rerun}

Run \texttt{python3\ rerun.py} in this workspace. It checks byte
preservation of the original input, verifies all authorized
primary-cache hashes, regenerates literal statement/proof units and the
nonnegative premise coverage index, performs the declared small finite
diagnostics, runs existing Pandoc and pdfLaTeX twice, and records
commands, environment, hashes and both complete pass logs.
\texttt{python3\ validate.py} checks only artifact consistency and the
backward note graph. These are documentary observations, never
mathematical correctness certificates.

Only the supplied input, the explicitly authorized read-only primary
files, public primary websites through the available web tool, and this
own fresh workspace were read. No other worker or model was used. A
direct public urllib fetch was denied by the network sandbox; native
public web retrieval supplied the additional primary text snapshots. One
mistakenly guessed arXiv identifier returned an unrelated
intersection-theory paper and one versioned old-style URL returned 404;
both were retained as retrieval diagnostics and supply no premise.
Temkin's introductory cofinality exercise is also not promoted to a
proof: the argument uses the stronger actual stable-modification theorem
instead.

The supplied 23 primary-file hashes match. The input and preserved copy
have 623435 bytes and SHA256
\texttt{e2f042118d966da4f8277f1f378116193539b60d039e4c0e3307525621cc719f}.
Source snapshots have their own exact text hashes; a hash of a
web-rendered snapshot is not falsely described as a downloaded PDF hash.
Existing tool versions and command outputs are retained. No supervised
compute service was supplied and no resource-heavy computation was
launched. The finite checks below are diagnostics on declared finite
rational grids, not evidence that the mathematical hypotheses are
fulfilled.

\texttt{status.json} has exactly
\texttt{id,status,exact\_result,proof\_file,concrete\_gaps}. The four
canonical notes contain this manuscript's backward-supported deductions
and explicit unresolved countable Runge/global C1 gate. Their supports
are only earlier locally registered note IDs; historical labels in the
input are not support handles. The full unchanged task text is appended
to the standalone paper by the rerun program. The failed first PDF pass
and its commands are retained separately; escaping an appendix backslash
corrected only typography, not the preserved task or mathematics.

\subsection{Primary premise index}\label{primary-premise-index}

\textbf{Premise 0.} BJ uses reduced possibly nonnormal projective
varieties over algebraically closed trivially valued fields. Definition
4.1 uses a decreasing net of generically finite FS functions on varying
rational bundles, whose numerical classes converge; bounded limits have
finite semiample approximants. Theorem 4.5 gives varying-class
decreasing closure. Source units: e001, e002.
\href{https://arxiv.org/pdf/1801.08229v4}{primary source}. Limit: Source
theorem used within its stated hypotheses; no independent proof
certificate.

\textbf{Premise 1.} For a projective finite-presentation flat family and
a finitely presented flat sheaf, Euler characteristics are locally
constant; generic flatness and generic geometric integrality permit
shrinking a finite-type integral field-of-definition base. These apply
separately to a fixed curve and fixed line bundle. Source units: e003,
e004, e005, e006.
\href{https://stacks.math.columbia.edu/tag/0B9T}{primary source};
\href{https://stacks.math.columbia.edu/tag/052A}{primary source};
\href{https://stacks.math.columbia.edu/tag/0574}{primary source};
\href{https://stacks.math.columbia.edu/tag/0553}{primary source}. Limit:
Source theorem used within its stated hypotheses; no independent proof
certificate.

\textbf{Premise 2.} BE Lemma 4.5 supplies finite presentation for
algebraic models over valuation rings. Proposition 5.13 extends a
projective morphism on a flat schematic graph model and identifies
pulled model metrics. Theorem 4.20 gives finite relative integral
closure for nontrivially valued discretely valued or algebraically
closed complete K, and reduced special fiber in the algebraically closed
case; Remark 4.21 excludes arbitrary densely valued K. Source units:
e007, e008, e009. \href{https://arxiv.org/pdf/1805.01016v3}{primary
source}. Limit: Source theorem used within its stated hypotheses; no
independent proof certificate.

\textbf{Premise 3.} Temkin defines multipointed curves as flat finitely
presented schemes with pure one-dimensional fibers. A multipointed curve
over integral qcqs S with semistable generic fiber has a stable
modification after generically etale alteration. Stable modification
morphisms are projective. Proposition 4.5.1 proves the valuation-ring
application for separably closed valued fields of finite height. Source
units: e010, e011, e012.
\href{https://arxiv.org/pdf/0707.3953v3}{primary source}. Limit: Source
theorem used within its stated hypotheses; no independent proof
certificate.

\textbf{Premise 4.} BPR identifies semistable vertex sets and semistable
formal models of a smooth proper curve over complete algebraically
closed nontrivially valued K; enlarging finite type II vertex sets gives
dominating refinements. This permits removal of loops while preserving
the already constructed graph morphism. Source units: e013, e014, e015.
\href{https://arxiv.org/pdf/1404.0279v1}{primary source};
\href{https://arxiv.org/pdf/1404.0279}{primary source}. Limit: Source
theorem used within its stated hypotheses; no independent proof
certificate.

\textbf{Premise 5.} BE continuous psh-regularizable metrics on
projective schemes define symmetric raw mixed Radon measures with
fundamental-cycle multiplicities and intersection mass. Real PL/model
metrics are dpsh, and reduced-special-fiber models give finite
Shilov-support coefficients equal to component intersections. FS metrics
admit smooth semipositive uniform regularization. Source units: e016,
e017. \href{https://arxiv.org/pdf/1805.01016v3}{primary source}. Limit:
Source theorem used within its stated hypotheses; no independent proof
certificate.

\textbf{Premise 6.} Wanner assumes K algebraically closed, complete and
nontrivially valued and works on smooth proper algebraic curves. Model
curvature agrees with Thuillier ddc in 3.27-3.28. Theorem 4.7, with its
proof, makes every finite continuous CLD-psh function on an ordinary
open Thuillier-subharmonic. Source units: e018, e019, e020.
\href{https://arxiv.org/pdf/1801.04713v3}{primary source}. Limit: Source
theorem used within its stated hypotheses; no independent proof
certificate.

\textbf{Premise 7.} Thuillier 3.3.7 gives local Green functions zero on
affinoid boundary with interior ddc=-delta and global proper-curve
functions with ddc(g\_xy)=delta\_y-delta\_x. Proposition 3.1.11 is the
subharmonic maximum principle, and 3.1.12 is its harmonic comparison
formulation. Source units: e021, e022.
\href{https://theses.hal.science/tel-00010990v1/document}{primary
source}. Limit: Source theorem used within its stated hypotheses; no
independent proof certificate.

\textbf{Premise 8.} Stacks 0AY8: a proper map to normal integral S from
reduced X, every generic point dominating, whose generic fiber has
global functions equal to the generic field, has geometrically connected
fibers and pushforward O=O. Stacks 0GIQ identifies finite normalization
of an integral geometrically unibranch scheme as a universal
homeomorphism. Source units: e023, e024.
\href{https://stacks.math.columbia.edu/tag/0AY8}{primary source};
\href{https://stacks.math.columbia.edu/tag/0GIQ}{primary source}. Limit:
Source theorem used within its stated hypotheses; no independent proof
certificate.

\textbf{Premise 9.} Stacks Bertini 0FD6 gives smooth general hyperplane
sections for the generated bundle and immersion hypotheses; 0FD9 gives
connected ample divisors when depth at closed points is at least two.
Characteristic-zero resolution by finite smooth-center blowups is used
for projective components, retaining projectivity. Source units: e025,
e026, e027, e028, e029.
\href{https://stacks.math.columbia.edu/tag/0FD6}{primary source};
\href{https://stacks.math.columbia.edu/tag/0FD9}{primary source};
\href{https://arxiv.org/pdf/alg-geom/9508005}{primary source};
\href{https://arxiv.org/abs/alg-geom/9508005}{primary source}. Limit:
Bierstone-Milman public alg-geom/9508005v2 is the introductory chapter,
with theorem statement and explanation; the complete chapter-4 proof was
not retrieved or reproved. This standard geometric theorem is an
explicit external input, not task-granted unconditional witness
construction.

\textbf{Premise 10.} A complete real-valued field norm extends uniquely
to algebraic extensions (Temkin Fact 2.2.2); an algebraic extension
embeds over its base into an algebraically closed field (Stacks 09GU).
Complete algebraically closed valued extensions are available for curve
tests. BJ Remark 1.9 supplies a continuous section of analytic
projection for non-Archimedean field extensions, hence surjectivity.
Source units: e030, e031, e032, e033.
\href{https://arxiv.org/pdf/1801.08229v4}{primary source};
\href{https://cml.centre-mersenne.org/item/10.1142/S1793744212500077.pdf}{primary
source};
\href{https://berkovich-2010.imj-prg.fr/lecture-notes/Temkin-lecture-notes.pdf}{primary
source}; \href{https://stacks.math.columbia.edu/tag/09GU}{primary
source}. Limit: Source theorem used within its stated hypotheses; no
independent proof certificate.

\textbf{Premise 11.} BJ 4.3 identifies constant potentials in nef
classes; Theorem 4.7 supplies finite maxima, sums, pullback and closure.
PSH topology is pointwise convergence on the divisorial valuations; sup
is continuous; 4.22 proves divisorial order. Corollary 12.18(ii), with
its 12.16 proof, supplies a decreasing sequence of finite rational FS
functions in varying ample bundles for an already genuine nef-class
potential, independently of target EP. No same-bundle regularization on
pi*L is asserted. Source units: e034, e035, e036, e037.
\href{https://arxiv.org/pdf/1801.08229v4}{primary source}. Limit: Source
theorem used within its stated hypotheses; no independent proof
certificate.

\textbf{Premise 12.} BJ 5.11 equates EP with compactness of normalized
PSH. Smooth connected characteristic-zero projective varieties have EP
for every nef class by 5.20, and compact normalized PSH by 5.21. No
birational onto statement is supplied before approaching target EP.
Source units: e038, e039.
\href{https://arxiv.org/pdf/1801.08229v4}{primary source}. Limit: Source
theorem used within its stated hypotheses; no independent proof
certificate.

\textbf{Premise 13.} Finite-presentation schemes, morphisms and finite
equality data descend along filtered affine qcqs inverse systems (Stacks
01ZM); finitely presented modules and vector bundles descend too (01ZR).
Relative ampleness of a line bundle is fpqc local on the base (0D3C).
Source units: e040, e041, e042.
\href{https://stacks.math.columbia.edu/tag/01ZM}{primary source};
\href{https://stacks.math.columbia.edu/tag/01ZR}{primary source};
\href{https://stacks.math.columbia.edu/tag/0D3C}{primary source}. Limit:
Source theorem used within its stated hypotheses; no independent proof
certificate.

\textbf{Premise 14.} Over a perfect field a normal variety is
geometrically normal (Stacks 0BXT). Noetherian flat local depth and
dimension formulas with regular fibers, and descent/ascent of
regularity, apply to finite-type models under characteristic-zero field
extension. BJ 1.13 and 1.15 characterize relative integral closure by
vertical R1 and special-fiber S1, preserving the specified nonnormal
generic fiber. Source units: e043, e044, e045, e046, e047.
\href{https://arxiv.org/pdf/1801.08229v4}{primary source};
\href{https://stacks.math.columbia.edu/tag/0BXT}{primary source};
\href{https://stacks.math.columbia.edu/tag/0337}{primary source};
\href{https://stacks.math.columbia.edu/tag/031E}{primary source};
\href{https://stacks.math.columbia.edu/tag/00OF}{primary source}. Limit:
Source theorem used within its stated hypotheses; no independent proof
certificate.

\textbf{Premise 15.} BJ 2.2 gives rational PL lattice density;
Definition 3.8 and Example 3.9 identify Hdom(c1(L)) with actual ample
models of L dominating the trivial model. Theorem 4.15(iii) supplies
decreasing Hdom majorant NETS for every ample-class genuine potential,
without needing EP. Source units: e048, e049, e050.
\href{https://arxiv.org/pdf/1801.08229v4}{primary source}. Limit: Source
theorem used within its stated hypotheses; no independent proof
certificate.

\textbf{Premise 16.} BE 8.3(iv) has raw complete-field pushforward
invariance with the analytic fundamental cycle; 8.4-8.5 uniquely extend
it under uniform metric convergence. BJ 4.16 provides uniform FS
approximation of continuous genuine ample-class potentials, and 7.23
identifies the actual continuous measures with CLD/BE on general reduced
projective X. Source units: e051, e052, e053.
\href{https://arxiv.org/pdf/1805.01016v3}{primary source};
\href{https://arxiv.org/pdf/1801.08229v4}{primary source}. Limit: Source
theorem used within its stated hypotheses; no independent proof
certificate.

\textbf{Premise 17.} BJ Theorem 13.1 orthogonality is conditional on EP
and uses differentiation of genuine envelope energy; Corollary 10.6 is
genuine finite-energy domination and does not require EP or an inverse
Dirac equation. Source units: e054, e055.
\href{https://arxiv.org/pdf/1801.08229v4}{primary source}. Limit: Source
theorem used within its stated hypotheses; no independent proof
certificate.

\textbf{Premise 18.} BJ 5.17 equates ample EP with continuity/genuine
membership of envelopes of all continuous obstacles; the envelope is
supnorm-Lipschitz. Lemma 5.9 and Theorem 4.5 give approaching-class
closure. These are applied after finite feasible-competitor bounds, not
assumed on a singular target. Source units: e056, e057.
\href{https://arxiv.org/pdf/1801.08229v4}{primary source}. Limit: Source
theorem used within its stated hypotheses; no independent proof
certificate.

\textbf{Premise 19.} BJ 5.13 gives birational pullback homeomorphism
only when the TARGET class is the limit of strictly larger target
classes with EP. It is used after target ample EP, with an explicit
conductor/Rees blowup presentation of normalization. Source units: e058.
\href{https://arxiv.org/pdf/1801.08229v4}{primary source}. Limit: Source
theorem used within its stated hypotheses; no independent proof
certificate.

\textbf{Premise 20.} BJ 7.2 gives normalization pullback equality for
mixed energy pairings. Theorem 7.18 and Proposition 7.19 define genuine
signed mixed finite-energy Radon measures, positivity, raw intersection
mass, actual central-cycle formula, decreasing continuity and
normalization pushforward identity (7.18), on the general reduced
projective variety. Source units: e059, e060.
\href{https://arxiv.org/pdf/1801.08229v4}{primary source}. Limit: Source
theorem used within its stated hypotheses; no independent proof
certificate.

\textbf{Premise 21.} BJ energy is the normalized mixed pairing with
denominator (n+1)V, extended by decreasing limits; E1 means energy
greater than minus infinity. The strong E1 topology refines the weak
topology by continuity of this energy, so equality of energy and common
valuations transports both. Source units: e061, e062.
\href{https://arxiv.org/pdf/1801.08229v4}{primary source}. Limit: Source
theorem used within its stated hypotheses; no independent proof
certificate.

\textbf{Premise 22.} BJ 9.1 defines dual energy for probability
measures. Theorem 9.20 on irreducible general X equates finite dual
energy with integrability of every genuine E1 potential. Proposition
9.19 supplies L1 continuity on bounded-energy weakly compact sets.
Source units: e063, e064.
\href{https://arxiv.org/pdf/1801.08229v4}{primary source}. Limit: Source
theorem used within its stated hypotheses; no independent proof
certificate.

\textbf{Premise 23.} BJ 8.3 identifies outer decreasing energy with
envelope energy under EP; Theorem 8.5 gives the uniform PL-test
derivative of genuine energy and Corollary 8.6 gives the continuous-test
derivative, on general varieties. Source units: e065.
\href{https://arxiv.org/pdf/1801.08229v4}{primary source}. Limit: Source
theorem used within its stated hypotheses; no independent proof
certificate.

\textbf{Premise 24.} BJ Theorem 12.8 equates ample EP with onto genuine
finite-energy MA on irreducible general X. Its proof uses actual
compactness, bounded-energy control, integral continuity, energy usc,
attainment and outer energy differentiation. No smooth-only corollary
replaces it. Source units: e066, e067.
\href{https://arxiv.org/pdf/1801.08229v4}{primary source}. Limit: Source
theorem used within its stated hypotheses; no independent proof
certificate.

\textbf{Premise 25.} BJ 7.40-7.44 gives genuine finite-energy
winning-locus locality and comparison; Corollary 10.4 gives uniqueness
modulo a real constant on irreducible X through actual E1 estimates and
divisorial order, without EP. Source units: e068, e069.
\href{https://arxiv.org/pdf/1801.08229v4}{primary source}. Limit: Source
theorem used within its stated hypotheses; no independent proof
certificate.

\textbf{Premise 26.} Wu 2609.27610v1 envelope proof expressly invokes
resolved-nef descent Proposition 2.24 and weak/global Theorem 6.5. Its
printed result is not an independent proof or review of C1 and is not
used as granted singular EP. Source units: e070.
\href{https://arxiv.org/pdf/2609.27610v1}{primary source}. Limit: Source
audit only. The theorem is not imported as established singular EP.

\subsection{Unchanged original mathematical task and full
family}\label{unchanged-original-mathematical-task-and-full-family}

This appendix reproduces the original task with typographic ASCII
normalization for pdfTeX. The authoritative exact text and JSON value
are preserved in \texttt{input-original.json} and
\texttt{task-original.json}. No mathematical condition has been
replaced.

Original P16.1: P16.1 (T5): Finite-energy equations on singular or
unibranch varieties Extend the finite-energy non-Archimedean MongeAmpere
operator and its existence-and- uniqueness theorem to singular or
unibranch varieties, replacing each use of smoothness in compactness,
envelopes, energy differentiation and comparison. State all retained
polarization and field hypotheses.

Complete unchanged mathematical family:

10 The general non-Archimedean MongeAmpere equation Active items: P16.1,
P16.2, P16.3, P16.4, P16.5, P16.6, P16.7. Status: Open in the residual
scope.

10.1 Atomic targets P16.1 (T5): Finite-energy equations on singular or
unibranch varieties Extend the finite-energy non-Archimedean MongeAmpere
operator and its existence-and- uniqueness theorem to singular or
unibranch varieties, replacing each use of smoothness in compactness,
envelopes, energy differentiation and comparison. State all retained
polarization and field hypotheses.

P16.2 (T5): Proper nonprojective algebraizable equation theories
Construct an appropriate replacement positivity datum and the associated
psh, finite-energy, comparison and equation theory on proper
nonprojective algebraizable spaces. A proper algebraic scheme with an
ample line bundle is already projective, so retaining that datum while
dropping projectivity is not a new case.

P16.3 (T5): Finite-energy equations in positive residue characteristic
Construct the required finite-energy operator and prove existence and
uniqueness when the original residue-characteristic restriction is
removed. Replace the resolution, model-regularity and envelope inputs
wherever they are used rather than merely changing the standing field
hypothesis.

P16.4 (T4): Finite-energy equations over arbitrary complete valued
fields Extend the operator and existence theorem to nondiscrete
valuations and arbitrary complete non-Archimedean fields outside the
established invariant-toric, skeleton-supported endpoint. A definition
of the finite-energy space alone does not establish the operator or the
equation.

P16.5 (T5): Finite-energy equations for nonample positivity classes
Extend the equation theory to big or pseudoeffective nonample classes,
specifying the psh cone, reference singularity type, energy and total
mass. At zero volume, the existing probability normalization must be
replaced rather than retained unchanged.

P16.6 (T4): Arbitrary-support finite-energy measures Prove the operator
and existence-and-uniqueness theorem for arbitrary finite-energy measure
support beyond the known toric skeletal regime. Measures must satisfy
the correct mass and energy conditions. Do not reduce arbitrary support
to a single skeleton or replace the genuine MongeAmpere measure by
another inverse measure construction.

P16.7 (T5): Finite-energy equations on nonalgebraizable analytic spaces
Construct the finite-energy operator and its existence-and-uniqueness
theory on genuinely nonalgebraizable analytic spaces, with an explicit
positivity datum and the necessary com- pactness, envelope and
comparison replacements.

\begin{verbatim}
                                          73
\end{verbatim}

10.2 What constitutes a solution For each axis, state the retained
hypotheses and construct the finite-energy operator together with the
corresponding existence-and-uniqueness theorem. Specify the measure mass
and finite-energy condition, uniqueness modulo constants and the
replacement inputs for com- pactness, envelopes, energy differentiation
and comparison. An endpoint solving a special toric measure class does
not remove a general axis.

10.3 Full statement and detailed state of the art Problem 10.1. Extend
the finite-energy existence and uniqueness theory for the NA Monge
Ampere equation beyond the currently proved combinations of smoothness,
projectivity, residue characteristic, valuation, positivity class,
support, and algebraizability hypotheses.

\begin{verbatim}
This is not one yes-or-no assertion. Positive characteristic, singular or unibranch varieties,
\end{verbatim}

nondiscrete valuations, big classes, arbitrary support, and
nonalgebraizable analytic spaces require different inputs. The known
variational proofs use compactness, envelope regularity,
differentiability of energy, and a comparison principle; deleting a
geometric hypothesis is valid only after replacing every place where it
enters. The relevant positive results and their distinct domains are
surveyed in Book Chapter 28, Book Chapter 29. The projectivity axis has
one logical constraint. A proper algebraic K-scheme endowed with an
ample line bundle is already projective over K {[}Stacksb, Tag 0B5N{]}.
Thus P16.2 cannot mean that one retains the BFJ line-bundle polarization
L while simply deleting the word projective. A nonprojective proper
algebraizable extension must first replace that positivity datumand then
construct the corresponding psh, finite-energy, and comparison
theoryrather than merely rerun the original argument. No canonical
replacement with the required equation theory is asserted here. One
exact source-stated axis is especially clean. PilleSchneider observes
that the finite- energy space can be defined over an arbitrary
non-Archimedean field, but that the extension of the MongeAmpere
operator is not known there in general; the discretely valued
characteristic- zero and trivially valued cases are the established
endpoints {[}PS23, Remark 4.4{]}. This should be solved before bundling
all remaining axes into a single theorem.

Proposition 10.1 (Arbitrary-field invariant toric skeletal endpoint).
Let K be a complete nontrivially valued non-Archimedean field, XSigma a
proper toric K-variety, and L an ample toric line bundle. If mu is a
positive Radon measure on XSigmaan , is supported on the canonical
skeleton NSigma , gives zero mass to NSigma \textbackslash NR , and has
total mass degL (XSigma ), then there is a toric continuous singular psh
metric on L whose MongeAmpere measure is mu. It is unique up to
multiplication of the metric by a positive constant.

\begin{verbatim}
Thus nondiscreteness and the arbitrary complete ground field cause no obstruction in
\end{verbatim}

this invariant toric, skeleton-supported range. Here continuous singular
has the sources extended-valued meaning: local weights are continuous
into R \{-infinity\}, as in {[}BGGJK25, 5.3.1{]}. It does not imply
finite real-valued weights or finite energy. This is not an arbitrary-
field finite-energy theorem in the BFJ sense: the result permits a
singular potential, requires torus invariance and skeletal support, and
does not construct the general extended Monge Ampere operator. In
particular, P16.4 and P16.6 retain the respectively nonskeletal and
noninvariant finite-energy residual scopes. There is also a distinct,
nonproper algebraic endpoint. Let n \textgreater{} d \textgreater= 1,
let X be a smooth projective Fano complex variety of dimension n, and
let D be a reduced simple-normal-crossing anticanonical divisor whose
dual complex is the standard (d-1)-simplex. For the affine Calabi Yau
complement X = X \textbackslash{} D, Wangs Theorem A constructs a
continuous non-Archimedean psh function on X an , over the trivially
valued field C, whose local MongeAmpere measure

\begin{verbatim}
                                           74
\end{verbatim}

is the canonical Lebesgue measure on the essential skeleton {[}Wan25{]}.
The source is a version- one preprint and uses a local
CLD/tropicalconvex construction. In particular, it neither constructs
the BFJ finite-energy operator nor treats arbitrary measures or
arbitrary affine varieties. Since X is affine and hence not proper, this
endpoint does not answer P16.2; it only records a special nonprojective
algebraic solution after the properness requirement is also dropped.
There is a separate finite-atomic inverse statement which must not be
called a finite-energy MongeAmpere theorem. Under its standing
assumptions that K is maximally complete and \textbar K x \textbar{} =
R\textgreater0 , Fang fixes an irreducible projective K-variety with an
ample line bundle and a finite set of graded birational points on its
affine cone. For every probability vector on that finite set, the source
constructs a Shilov-finite, possibly singular metric having the
prescribed equidistribution measure {[}Fan25, Proposition 4.8{]}. The
source itself says that comparing this Banach-algebra measure with the
genuine MongeAmpere measure requires the separate con- tinuity property
of an envelope metric. Thus it supplies neither a continuous
semipositive solution nor the BFJ finite-energy operator for the
prescribed atomic measure, and does not narrow P16.6. There is also a
resolved negative regularity chronology. The first
skeletal-factorization question of GublerJellKunnemannMartin already
fails on their toric surface {[}GJKM19, Appendix A{]}. BoucksomJonsson
then construct a finite atomic MongeAmpere measure sup- ported on a dual
complex whose solution is invariant under no retraction attached to any
SNC model; they also construct a divisorial Green function that is not
Q-PL {[}BJ24b, Theo- rems A and B{]}. Hence support on a dual complex
does not imply factorization through some refined skeleton. This is a
theorem-level obstruction to a proof strategy, not an open regularity
expectation. The Rethlas programme treats each axis separately and first
searches the post-2020 liter- ature for an existing theorem. Alterations
and universal homeomorphisms can address some algebraic descent
questions, but they do not by themselves supply multiplier ideals or
dif- ferentiability in positive characteristic. The unrestricted
formulation is open; several strictly delimited axes are already known.

\begin{verbatim}
                                           75
\end{verbatim}

\textbf{Exact completion criterion (typography normalized).} Establish
the genuine finite-energy operator, exact measure mass and energy
condition, existence, uniqueness modulo constants, and every
compactness/envelope/differentiation/comparison replacement for the
singular/geometrically unibranch projective domain with explicitly
retained field and polarization hypotheses. Verify the new envelope
source and the actual general-variety global theory. No change to the
original mathematical target is permitted.

\end{document}
